\section{Metodologi}

\subsection{Setup Eksperimen}
\begin{table}[h]
\centering
\caption{Setup Eksperimen}
\scriptsize
\begin{tabular}{|l|l|}
\hline
\textbf{Parameter} & \textbf{Value} \\
\hline
\multicolumn{2}{|c|}{\textbf{Training Hyperparameters}} \\
\hline
Epochs & 15 \\
Batch Size & 32 \\
Learning Rate & $1 \times 10^{-4}$ \\
Optimizer & AdamW \\
Weight Decay & $1 \times 10^{-4}$ \\
Loss Function & CrossEntropyLoss \\
\hline
\multicolumn{2}{|c|}{\textbf{Data Preprocessing}} \\
\hline
Image Size & 224 $\times$ 224 \\
RandomHorizontalFlip & p = 0.5 \\
ColorJitter & brightness/contrast/saturation = 0.2 \\
Normalization Mean & [0.485, 0.456, 0.406] \\
Normalization Std & [0.229, 0.224, 0.225] \\
\hline
\multicolumn{2}{|c|}{\textbf{Hardware \& Reproducibility}} \\
\hline
Device & CUDA GPU \\
Precision & Mixed (fp16/bf16) \\
Random Seed & 42 \\
\hline
\end{tabular}
\end{table}

\subsection{Arsitektur Model}

\begin{table}[h]
\centering
\caption{Arsitektur Model yang Diuji}
\scriptsize
\begin{tabular}{|p{2.5cm}|>{\raggedright\arraybackslash}p{11cm}|}
\hline
\textbf{Model} & \textbf{Deskripsi Arsitektur} \\
\hline
\textbf{Plain-34} & Jaringan dasar 34 layer tanpa skip connection. Setiap blok berisi dua Conv3x3 dengan BN dan ReLU. Karena tidak ada shortcut, model sulit dilatih saat semakin dalam \\
\hline
\textbf{ResNet-34} & Sama dengan Plain-34, tapi setiap blok punya skip connection ($x+F(x)$). Fitur utama adalah shortcut yang menjaga aliran informasi dan gradient, sehingga training lebih stabil \\
\hline
\textbf{Pre-activation ResNet-34} & Modifikasi ResNet-34 dengan urutan BN-ReLU sebelum Conv, serta tanpa aktivasi setelah penjumlahan skip \\
\hline
\textbf{SE-ResNet-34} & ResNet-34 yang ditambah modul Squeeze-and-Excitation. Modul ini memberi perhatian (attention) pada channel penting, memperkuat fitur relevan, dan hanya menambah parameter sedikit ($<1\%$). \\
\hline
\textbf{SE-Preact ResNet-34} & Kombinasi pre-activation block dengan SE module. Tujuannya menggabungkan jalur gradient yang lebih bersih dengan channel attention, meski pada depth 34 layer keuntungannya terbatas. \\
\hline
\end{tabular}
\end{table}
